\chapter*{用語索引}
\addcontentsline{toc}{chapter}{用語索引}
\markboth{用語索引}{}
主な用語と,定義・説明を置いた箇所を示す.頁番号と用語名から本文へ移動できる.
\begingroup
\renewcommand{\arraystretch}{1.2}
\begin{longtable}{@{}p{90mm}p{35mm}r@{}}
\toprule
用語 & 主な参照先 & 頁\\
\midrule
\endfirsthead
\toprule
用語 & 主な参照先 & 頁\\
\midrule
\endhead
\hyperref[節:一次近似]{一次近似} & \sectionref{節:一次近似} & \pageref{節:一次近似}\\
\hyperref[節:運動と積分]{位置の変化と道のり} & \sectionref{節:運動と積分} & \pageref{節:運動と積分}\\
\hyperref[節:逆媒介陰微分]{陰関数の微分} & \sectionref{節:逆媒介陰微分} & \pageref{節:逆媒介陰微分}\\
\hyperref[節:ウォリス公式]{ウォリスの公式} & \sectionref{節:ウォリス公式} & \pageref{節:ウォリス公式}\\
\hyperref[式:円筒殻の体積]{円筒殻} & 式\eqref{式:円筒殻の体積} & \pageref{式:円筒殻の体積}\\
\hyperref[節:回転体の体積]{回転体の体積} & \sectionref{節:回転体の体積} & \pageref{節:回転体の体積}\\
\hyperref[式:回転面の面積]{回転面の面積} & 式\eqref{式:回転面の面積} & \pageref{式:回転面の面積}\\
\hyperref[節:等比級数]{級数と部分和} & \sectionref{節:等比級数} & \pageref{節:等比級数}\\
\hyperref[節:曲線の弧長]{曲線の長さ} & \sectionref{節:曲線の弧長} & \pageref{節:曲線の弧長}\\
\hyperref[節:奇偶の対称性]{奇関数・偶関数} & \sectionref{節:奇偶の対称性} & \pageref{節:奇偶の対称性}\\
\hyperref[節:逆媒介陰微分]{逆関数の微分} & \sectionref{節:逆媒介陰微分} & \pageref{節:逆媒介陰微分}\\
\hyperref[節:逆三角関数]{逆三角関数} & \sectionref{節:逆三角関数} & \pageref{節:逆三角関数}\\
\hyperref[節:和の区分求積]{区分求積法} & \sectionref{節:和の区分求積} & \pageref{節:和の区分求積}\\
\hyperref[節:原始関数基本定理]{原始関数} & \sectionref{節:原始関数基本定理} & \pageref{節:原始関数基本定理}\\
\hyperref[節:非初等原始関数]{原始関数の存在と表示} & \sectionref{節:非初等原始関数} & \pageref{節:非初等原始関数}\\
\hyperref[節:円の有理パラメータ]{原始ピタゴラス数} & \sectionref{節:円の有理パラメータ} & \pageref{節:円の有理パラメータ}\\
\hyperref[節:広義積分の定義]{広義積分} & \sectionref{節:広義積分の定義} & \pageref{節:広義積分の定義}\\
\hyperref[節:積分の不等式]{コーシー・シュワルツの不等式} & \sectionref{節:積分の不等式} & \pageref{節:積分の不等式}\\
\hyperref[節:微分法則]{合成関数の微分} & \sectionref{節:微分法則} & \pageref{節:微分法則}\\
\hyperref[節:周期関数の積分]{周期関数の積分} & \sectionref{節:周期関数の積分} & \pageref{節:周期関数の積分}\\
\hyperref[節:数列極限]{収束と発散} & \sectionref{節:数列極限} & \pageref{節:数列極限}\\
\hyperref[節:数値積分]{シンプソン公式} & \sectionref{節:数値積分} & \pageref{節:数値積分}\\
\hyperref[節:面積重心と線重心]{重心（面積・線）} & \sectionref{節:面積重心と線重心} & \pageref{節:面積重心と線重心}\\
\hyperref[節:部分積分の循環]{循環する部分積分} & \sectionref{節:部分積分の循環} & \pageref{節:部分積分の循環}\\
\hyperref[節:相反三角関数]{正割・余割・余接} & \sectionref{節:相反三角関数} & \pageref{節:相反三角関数}\\
\hyperref[節:定積分性質]{積分の平均値の定理} & \sectionref{節:定積分性質} & \pageref{節:定積分性質}\\
\hyperref[節:変数上下端]{積分区間の端が変数の場合} & \sectionref{節:変数上下端} & \pageref{節:変数上下端}\\
\hyperref[節:原始関数基本定理]{積分定数} & \sectionref{節:原始関数基本定理} & \pageref{節:原始関数基本定理}\\
\hyperref[節:積分方程式]{積分方程式} & \sectionref{節:積分方程式} & \pageref{節:積分方程式}\\
\hyperref[節:積分漸化式]{積分漸化式} & \sectionref{節:積分漸化式} & \pageref{節:積分漸化式}\\
\hyperref[節:双曲線関数]{双曲線関数} & \sectionref{節:双曲線関数} & \pageref{節:双曲線関数}\\
\hyperref[節:数値積分]{台形公式} & \sectionref{節:数値積分} & \pageref{節:数値積分}\\
\hyperref[式:中央二項係数の極限]{中央二項係数} & 式\eqref{式:中央二項係数の極限} & \pageref{式:中央二項係数の極限}\\
\hyperref[章:置換積分]{置換積分} & \chapterref{章:置換積分} & \pageref{章:置換積分}\\
\hyperref[章:定積分の対称性]{定積分の対称性} & \chapterref{章:定積分の対称性} & \pageref{章:定積分の対称性}\\
\hyperref[節:有限テイラー]{テイラー公式と剰余項} & \sectionref{節:有限テイラー} & \pageref{節:有限テイラー}\\
\hyperref[節:等比級数]{等比級数} & \sectionref{節:等比級数} & \pageref{節:等比級数}\\
\hyperref[節:微分定義]{導関数} & \sectionref{節:微分定義} & \pageref{節:微分定義}\\
\hyperref[節:取り尽くし法]{取り尽くし法} & \sectionref{節:取り尽くし法} & \pageref{節:取り尽くし法}\\
\hyperref[節:ニュートン法]{ニュートン法} & \sectionref{節:ニュートン法} & \pageref{節:ニュートン法}\\
\hyperref[節:半角置換]{半角置換} & \sectionref{節:半角置換} & \pageref{節:半角置換}\\
\hyperref[節:曲線の弧長]{媒介曲線の面積・長さ} & \sectionref{節:曲線の弧長} & \pageref{節:曲線の弧長}\\
\hyperref[節:パップスの定理]{パップス＝ギュルダンの定理} & \sectionref{節:パップスの定理} & \pageref{節:パップスの定理}\\
\hyperref[節:基本定理第一]{微積分の基本定理} & \sectionref{節:基本定理第一} & \pageref{節:基本定理第一}\\
\hyperref[節:微分定義]{微分係数} & \sectionref{節:微分定義} & \pageref{節:微分定義}\\
\hyperref[節:微分方程式]{微分方程式} & \sectionref{節:微分方程式} & \pageref{節:微分方程式}\\
\hyperref[節:原始関数基本定理]{不定積分} & \sectionref{節:原始関数基本定理} & \pageref{節:原始関数基本定理}\\
\hyperref[節:部分分数]{部分分数} & \sectionref{節:部分分数} & \pageref{節:部分分数}\\
\hyperref[章:部分積分]{部分積分} & \chapterref{章:部分積分} & \pageref{章:部分積分}\\
\hyperref[節:部分積分の反復]{部分積分の反復} & \sectionref{節:部分積分の反復} & \pageref{節:部分積分の反復}\\
\hyperref[節:平均値定理]{平均値の定理（微分）} & \sectionref{節:平均値定理} & \pageref{節:平均値定理}\\
\hyperref[節:関数の平均]{平均値（関数）} & \sectionref{節:関数の平均} & \pageref{節:関数の平均}\\
\hyperref[節:広義積分の判定]{べき関数による収束判定} & \sectionref{節:広義積分の判定} & \pageref{節:広義積分の判定}\\
\hyperref[節:ベータ積分]{ベータ関数} & \sectionref{節:ベータ積分} & \pageref{節:ベータ積分}\\
\hyperref[節:曲線間の面積]{面積（曲線で囲まれた領域）} & \sectionref{節:曲線間の面積} & \pageref{節:曲線間の面積}\\
\hyperref[節:リーマン和]{リーマン和} & \sectionref{節:リーマン和} & \pageref{節:リーマン和}\\
\hyperref[節:関数極限連続]{連続性} & \sectionref{節:関数極限連続} & \pageref{節:関数極限連続}\\
\hyperref[用語:連続微分可能]{連続微分可能} & \sectionref{節:合成形の置換} & \pageref{用語:連続微分可能}\\
\hyperref[節:ロピタル]{ロピタルの定理} & \sectionref{節:ロピタル} & \pageref{節:ロピタル}\\
\hyperref[節:半角置換]{ワイエルシュトラス置換} & \sectionref{節:半角置換} & \pageref{節:半角置換}\\
\hyperref[節:相補重み]{和が一になる関数} & \sectionref{節:相補重み} & \pageref{節:相補重み}\\
\hyperref[節:和と積分の比較]{和と積分の比較} & \sectionref{節:和と積分の比較} & \pageref{節:和と積分の比較}\\
\bottomrule
\end{longtable}
\endgroup
